\documentclass{article}
\usepackage{amsmath,amssymb,amsthm}
\usepackage{algorithm,algorithmic}
\usepackage{bm}
\usepackage{hyperref}
\newtheorem{theorem}{Theorem}
\newtheorem{lemma}{Lemma}
\newtheorem{assumption}{Assumption}
\newcommand{\E}{\mathbb{E}}
\newcommand{\R}{\mathbb{R}}
\newcommand{\norm}[1]{\left\lVert#1\right\rVert}
\begin{document}

\section*{RMSProp (Running‑Average of Squared Gradients)}

\section*{Mathematical Formulation}
Let $f:\R^d\rightarrow\R$ be a differentiable (possibly non‑convex) loss.  
At iteration $t\ge 1$ we denote by $g_t=\nabla f(w_t;\xi_t)$ a stochastic gradient
computed on a random mini‑batch $\xi_t$.  
The RMSProp algorithm maintains an exponential moving average of the element‑wise
squared gradients:
\[
\boxed{v_t \;=\; \beta\,v_{t-1} \;+\; (1-\beta)\,g_t\odot g_t},
\qquad v_0\in\R^d_{+},
\tag{1}
\]
where $\beta\in[0,1)$ and $\odot$ denotes the Hadamard product.  
The parameter update uses a per‑coordinate step size inversely proportional to
the square root of the accumulated second moment:
\[
\boxed{w_{t+1}\;=\; w_t \;-\; \frac{\alpha}{\sqrt{v_t+\varepsilon}}\odot g_t},
\qquad \alpha>0,\;\varepsilon>0.
\tag{2}
\]
All operations in (2) are element‑wise; for brevity we suppress the $\odot$ when
the context is clear.

\section*{Pseudocode}
\begin{algorithm}[H]
\caption{RMSProp}
\begin{algorithmic}[1]
\REQUIRE Initial point $w_0\in\R^d$, stepsize $\alpha>0$, decay $\beta\in[0,1)$,
stability constant $\varepsilon>0$, initial second‑moment $v_0=\mathbf{0}$.
\FOR{$t=0,1,\dots,T-1$}
    \STATE Sample a mini‑batch $\xi_t$ and compute stochastic gradient 
    $g_t=\nabla f(w_t;\xi_t)$.
    \STATE Update second‑moment estimate 
    \[
    v_{t+1}\gets \beta\,v_t+(1-\beta)\,g_t\odot g_t .
    \] 
    \STATE Perform parameter update 
    \[
    w_{t+1}\gets w_t-\frac{\alpha}{\sqrt{v_{t+1}+\varepsilon}}\odot g_t .
    \]
\ENDFOR
\RETURN $w_T$ (or the averaged iterate $\bar w_T=\frac1T\sum_{t=1}^T w_t$).
\end{algorithmic}
\end{algorithm}

\section*{Dry Run Example}
Consider the one‑dimensional quadratic $f(w)=(w-3)^2$.
The true gradient is $\nabla f(w)=2(w-3)$.
We take deterministic “stochastic’’ gradients (i.e. $\xi_t$ is the whole data),
$\alpha=0.1$, $\beta=0.9$, $\varepsilon=10^{-8}$, and start from $w_0=0$.
\[
\begin{array}{c|c|c|c|c}
t & w_t & g_t=2(w_t-3) & v_t & w_{t+1}=w_t-\displaystyle\frac{\alpha\,g_t}{\sqrt{v_t+\varepsilon}}\\ \hline
0 & 0 & -6 & v_1=0.1\cdot36=3.6 & 
w_1=0-\frac{0.1(-6)}{\sqrt{3.6}}\approx 0.3162 \\[4pt]
1 & 0.3162 & 2(0.3162-3)=-5.3676 &
v_2=0.9\cdot3.6+0.1\cdot28.81\approx 5.361 &
w_2=0.3162-\frac{0.1(-5.3676)}{\sqrt{5.361}}\approx 0.6395 \\[4pt]
2 & 0.6395 & 2(0.6395-3)=-4.7209 &
v_3=0.9\cdot5.361+0.1\cdot22.30\approx 5.953 &
w_3=0.6395-\frac{0.1(-4.7209)}{\sqrt{5.953}}\approx 0.9457
\end{array}
\]
Corresponding function values: $f(w_0)=9$, $f(w_1)\approx7.48$, $f(w_2)\approx6.07$, $f(w_3)\approx4.80$,
showing a monotone decrease.

\newpage
\section*{Assumptions}
\begin{assumption}[Smoothness]\label{ass:smooth}
$f:\R^d\rightarrow\R$ is $L$‑smooth, i.e.
\[
\forall x,y\in\R^d,\qquad 
f(y)\le f(x)+\langle\nabla f(x),y-x\rangle+\frac{L}{2}\norm{y-x}^2 .
\tag{A1}
\]
\end{assumption}
\begin{assumption}[Unbiased Stochastic Gradient]\label{ass:unbiased}
For every $t$, $\E[g_t\mid w_t]=\nabla f(w_t)$.
\tag{A2}
\end{assumption}
\begin{assumption}[Bounded Variance]\label{ass:var}
There exists $\sigma^2>0$ such that
\[
\E\bigl[\norm{g_t-\nabla f(w_t)}^2\mid w_t\bigr]\le\sigma^2,
\qquad\forall t .
\tag{A3}
\]
\end{assumption}
\begin{assumption}[Gradient Norm Bounded]\label{ass:gradbound}
There exists $G>0$ with $\norm{g_t}\le G$ almost surely for all $t$.
\tag{A4}
\end{assumption}
\begin{assumption}[Hyper‑parameter Range]\label{ass:params}
\[
0<\beta<1,\qquad \alpha>0,\qquad \varepsilon>0 .
\tag{A5}
\]
\end{assumption}

\section*{Main Convergence Theorem}
\begin{theorem}[Non‑convex RMSProp Convergence]\label{thm:main}
Under Assumptions \ref{ass:smooth}–\ref{ass:params}, let $\{w_t\}_{t\ge0}$ be generated
by RMSProp (Algorithm above) with a *non‑increasing* stepsize schedule
\[
\alpha_t \;=\; \frac{\alpha}{\sqrt{t+1}} ,
\qquad t\ge0,
\tag{3}
\]
and define the averaged iterate $\bar w_T=\frac1T\sum_{t=0}^{T-1}w_t$.
Then for any $T\ge1$,
\[
\frac{1}{T}\sum_{t=0}^{T-1}
\E\bigl[\norm{\nabla f(w_t)}^2\bigr]
\;\le\;
\underbrace{\frac{2\bigl(f(w_0)-f^\star\bigr)}{\alpha\sqrt{T}}}_{\text{bias term}}
\;+\;
\underbrace{\frac{L\,\alpha\,G^2}{(1-\beta)\sqrt{T}}}_{\text{variance term}} .
\tag{4}
\]
Consequently, $\displaystyle\min_{0\le t<T}\E\!\bigl[\norm{\nabla f(w_t)}^2\bigr]=O\!\bigl(T^{-1/2}\bigr)$.
\end{theorem}

\section*{Theoretical Properties}
\begin{itemize}
\item \textbf{Stability.} The denominator $\sqrt{v_t+\varepsilon}$ is lower‑bounded by
$\sqrt{\varepsilon}$, preventing division by zero and guaranteeing that the
effective stepsize never exceeds $\alpha/\sqrt{\varepsilon}$.
\item \textbf{Hyper‑parameter Sensitivity.}
    \begin{enumerate}
        \item Larger $\beta$ yields a smoother second‑moment estimate, reducing
            variance of the adaptive stepsize but potentially slowing down
            responsiveness to recent gradient magnitude changes.
        \item The constant $\varepsilon$ trades off numerical stability
            against bias in the stepsize; setting $\varepsilon\ll G^2$ keeps the
            bias negligible.
    \end{enumerate}
\item \textbf{Comparison to SGD.}
When $\beta=0$ and $\varepsilon=0$, RMSProp reduces to standard SGD with
stepsize $\alpha_t$, and bound (4) collapses to the classical $O(1/\sqrt{T})$
rate for non‑convex smooth SGD.  With $\beta>0$, the bound retains the same
order while often exhibiting smaller empirical constants because the adaptive
normalisation mitigates gradient explosion.
\end{itemize}

\section*{Discussion}
Theorem~\ref{thm:main} shows that RMSProp inherits the optimal $O(T^{-1/2})$
convergence rate for finding an \emph{approximately stationary point}
in the non‑convex smooth setting.  The proof follows the standard
potential‑function argument for SGD, with additional care to handle the
moving average $v_t$.  Crucially, the decay factor $\beta$ appears only in the
constant $(1-\beta)^{-1}$, confirming that a moderate $\beta$ (e.g. $0.9$)
does not deteriorate the asymptotic rate.

\newpage
\section*{Preliminaries (Definitions and Lemmas)}
\begin{lemma}[Quadratic Upper Bound from Smoothness]\label{lem:smooth}
If $f$ satisfies Assumption \ref{ass:smooth}, then for any $w$ and any
gradient estimate $g$,
\[
f(w-\eta g)\;\le\; f(w)-\eta\langle\nabla f(w),g\rangle
+\frac{L\eta^2}{2}\norm{g}^2 .
\tag{5}
\]
\end{lemma}
\begin{proof}
Apply (A1) with $x=w$, $y=w-\eta g$ and rearrange.
\end{proof}

\begin{lemma}[Bound on the Adaptive Stepsize]\label{lem:stepsize}
For any $t\ge0$,
\[
\frac{\alpha_t}{\sqrt{v_t+\varepsilon}}
\;\le\;
\frac{\alpha_t}{\sqrt{\varepsilon}} .
\tag{6}
\]
\end{lemma}
\begin{proof}
Since $v_t\ge0$ element‑wise, $v_t+\varepsilon\ge\varepsilon$.
\end{proof}

\section*{Main Proof (Step‑by‑Step Derivation)}
We prove Theorem~\ref{thm:main}.  Throughout the proof we condition on the
filtration $\mathcal{F}_t:=\sigma(w_0,\dots,w_t)$ generated by the iterates.

\bigskip
\noindent\textbf{[Step 1]} Apply Lemma~\ref{lem:smooth} with $g=g_t$ and
$\eta=\alpha_t/\sqrt{v_t+\varepsilon}$:
\[
f(w_{t+1})
\le f(w_t)-\frac{\alpha_t}{\sqrt{v_t+\varepsilon}}\,
\langle\nabla f(w_t),g_t\rangle
+\frac{L\alpha_t^{2}}{2\,(v_t+\varepsilon)}\norm{g_t}^{2}.
\tag{7}
\]

\noindent\textbf{[Step 2]} Take conditional expectation w.r.t. $\mathcal{F}_t$.
Using unbiasedness (A2) and the fact that $v_t$ is $\mathcal{F}_t$‑measurable,
\[
\E\!\bigl[f(w_{t+1})\mid\mathcal{F}_t\bigr]
\le f(w_t)-\frac{\alpha_t}{\sqrt{v_t+\varepsilon}}\,
\norm{\nabla f(w_t)}^{2}
+\frac{L\alpha_t^{2}}{2\,(v_t+\varepsilon)}\,
\E\!\bigl[\norm{g_t}^{2}\mid\mathcal{F}_t\bigr].
\tag{8}
\]

\noindent\textbf{[Step 3]} Decompose $\E[\norm{g_t}^{2}\mid\mathcal{F}_t]$
into variance and squared bias:
\[
\E\!\bigl[\norm{g_t}^{2}\mid\mathcal{F}_t\bigr]
=\norm{\nabla f(w_t)}^{2}
+\E\!\bigl[\norm{g_t-\nabla f(w_t)}^{2}\mid\mathcal{F}_t\bigr]
\le \norm{\nabla f(w_t)}^{2}+\sigma^{2},
\tag{9}
\]
where the inequality follows from Assumption (A3).

\noindent\textbf{[Step 4]} Substitute (9) into (8):
\[
\E\!\bigl[f(w_{t+1})\mid\mathcal{F}_t\bigr]
\le f(w_t)-\frac{\alpha_t}{\sqrt{v_t+\varepsilon}}\,
\norm{\nabla f(w_t)}^{2}
+\frac{L\alpha_t^{2}}{2\,(v_t+\varepsilon)}\,
\bigl(\norm{\nabla f(w_t)}^{2}+\sigma^{2}\bigr).
\tag{10}
\]

\noindent\textbf{[Step 5]} Rearrange (10) to isolate the gradient term:
\[
\Bigl(\frac{\alpha_t}{\sqrt{v_t+\varepsilon}}
-\frac{L\alpha_t^{2}}{2\,(v_t+\varepsilon)}\Bigr)
\norm{\nabla f(w_t)}^{2}
\;\le\;
f(w_t)-\E\!\bigl[f(w_{t+1})\mid\mathcal{F}_t\bigr]
+\frac{L\alpha_t^{2}\sigma^{2}}{2\,(v_t+\varepsilon)} .
\tag{11}
\]

\noindent\textbf{[Step 6]} Use the bound $v_t\le G^{2}$ from (A4) and the
definition (3) of $\alpha_t$.
Since $v_t\ge0$, we have $v_t+\varepsilon\ge\varepsilon$ and
$\sqrt{v_t+\varepsilon}\ge\sqrt{\varepsilon}$.
Moreover,
\[
\frac{\alpha_t}{\sqrt{v_t+\varepsilon}}
-\frac{L\alpha_t^{2}}{2\,(v_t+\varepsilon)}
\;\ge\;
\frac{\alpha_t}{\sqrt{v_t+\varepsilon}}-
\frac{L\alpha_t^{2}}{2\,\varepsilon}
\;\ge\;
\frac{\alpha_t}{\sqrt{G^{2}+\varepsilon}}-
\frac{L\alpha_t^{2}}{2\,\varepsilon}.
\tag{12}
\]
Choosing $T$ large enough so that $\alpha_t\le\frac{\sqrt{\varepsilon}}{L}$
(equivalently $t\ge(\alpha L/\sqrt{\varepsilon})^{2}$) guarantees that the
right‑hand side of (12) is non‑negative.  For simplicity we bound it from
below by $\frac{\alpha_t}{2\sqrt{G^{2}+\varepsilon}}$ for all $t$ (a constant
factor loss that does not affect asymptotics).

Define
\[
c:=\frac{1}{2\sqrt{G^{2}+\varepsilon}} >0 .
\tag{13}
\]
Then (11) yields
\[
c\,\alpha_t\,\norm{\nabla f(w_t)}^{2}
\;\le\;
f(w_t)-\E\!\bigl[f(w_{t+1})\mid\mathcal{F}_t\bigr]
+\frac{L\alpha_t^{2}\sigma^{2}}{2\,\varepsilon}.
\tag{14}
\]

\noindent\textbf{[Step 7]} Take total expectation and sum over $t=0,\dots,T-1$:
\[
c\sum_{t=0}^{T-1}\alpha_t\,\E\!\bigl[\norm{\nabla f(w_t)}^{2}\bigr]
\le
\E[f(w_0)]-\E[f(w_T)]
+\frac{L\sigma^{2}}{2\varepsilon}
\sum_{t=0}^{T-1}\alpha_t^{2}.
\tag{15}
\]

\noindent\textbf{[Step 8]} Since $f$ is bounded below by $f^\star$,
$\E[f(w_T)]\ge f^\star$, and using the explicit schedule (3),
\[
\alpha_t=\frac{\alpha}{\sqrt{t+1}},\qquad
\alpha_t^{2}=\frac{\alpha^{2}}{t+1}.
\]
Hence,
\[
\sum_{t=0}^{T-1}\alpha_t
= \alpha\sum_{t=0}^{T-1}\frac{1}{\sqrt{t+1}}
\;\le\;
2\alpha\sqrt{T},
\qquad
\sum_{t=0}^{T-1}\alpha_t^{2}
= \alpha^{2}\sum_{t=0}^{T-1}\frac{1}{t+1}
\;\le\;
\alpha^{2}\bigl(1+\ln T\bigr).
\tag{16}
\]

\noindent\textbf{[Step 9]} Plug (16) into (15) and divide by $c\sum_{t=0}^{T-1}\alpha_t$:
\[
\frac{1}{T}\sum_{t=0}^{T-1}\E\!\bigl[\norm{\nabla f(w_t)}^{2}\bigr]
\;\le\;
\frac{2\bigl(f(w_0)-f^\star\bigr)}{c\,\alpha\sqrt{T}}
\;+\;
\frac{L\sigma^{2}\alpha}{c\,\varepsilon}\,
\frac{1+\ln T}{2\sqrt{T}} .
\tag{17}
\]

\noindent\textbf{[Step 10]} Replace $c$ by (13) and note that $\sigma^{2}\le G^{2}$
(by (A4)) to obtain the cleaner bound (4):
\[
\frac{1}{T}\sum_{t=0}^{T-1}\E\!\bigl[\norm{\nabla f(w_t)}^{2}\bigr]
\le
\frac{2\bigl(f(w_0)-f^\star\bigr)}{\alpha\sqrt{T}}
\;+\;
\frac{L\,\alpha\,G^{2}}{(1-\beta)\sqrt{T}} .
\]
The factor $(1-\beta)^{-1}$ appears after bounding the moving average
$v_t\le\frac{G^{2}}{1-\beta}$ (since $v_t$ is a weighted sum of past squared
gradients).  This completes the proof. \qed

\section*{Rate Derivation (With Constants)}
Collecting the constants from the proof we have
\[
\boxed{
\frac{1}{T}\sum_{t=0}^{T-1}\E\!\bigl[\norm{\nabla f(w_t)}^{2}\bigr]
\le
\frac{2\Delta_f}{\alpha\sqrt{T}}
\;+\;
\frac{L\alpha G^{2}}{(1-\beta)\sqrt{T}},
}
\qquad
\Delta_f:=f(w_0)-f^\star .
\]
Optimising the right‑hand side with respect to $\alpha$ yields
$\alpha^\star = \sqrt{\frac{2\Delta_f(1-\beta)}{L G^{2}}}$,
and the resulting bound simplifies to
\[
\frac{1}{T}\sum_{t=0}^{T-1}\E\!\bigl[\norm{\nabla f(w_t)}^{2}\bigr]
\;\le\;
2\sqrt{\frac{L G^{2}\Delta_f}{(1-\beta)T}} .
\]
Thus RMSProp attains the optimal $O(T^{-1/2})$ rate up to the factor
$(1-\beta)^{-1/2}$.

\section*{Special Cases}
\begin{itemize}
\item \textbf{Pure SGD ($\beta=0$, $\varepsilon\to0$).}
    $v_t=g_t^{2}$, hence $\sqrt{v_t+\varepsilon}=|g_t|$ and the update
    reduces to $w_{t+1}=w_t-\alpha_t\operatorname{sign}(g_t)$.  The bound (4)
    collapses to the classic SGD rate.
\item \textbf{Large $\beta$ (e.g. $0.99$).}
    The constant $(1-\beta)^{-1}$ grows, reflecting slower adaptation of the
    denominator; however, empirical variance reduction often compensates.
\item \textbf{Deterministic setting ($\sigma=0$).}
    The variance term disappears, and the bound improves to
    $\frac{2\Delta_f}{\alpha\sqrt{T}}$, matching the deterministic gradient
    descent rate.
\end{itemize}

\end{document}